\SourceLeaf{29}\noindent also ist $\varphi t$ notwendig durch $kk$ teilbar, und die Koeffizienten in $Ft'$ sind durch
$\dfrac{\varphi t}{kk}$ teilbar.

Sei $k$ der größte gemeinsame Teiler der Koeffizienten in $Ft$ und sei $Ft = ku$. Ist
$l$ der größte gemeinsame Teiler der Koeffizienten in $Fu$, so gilt
\[
Fu = lt, \qquad \varphi t = kkl, \qquad \varphi u = kll, \qquad tu = kl.
\]

Ferner sind die Koeffizienten von $\nu\nu$ in $t$, $u$ jeweils teilerfremd zu $k$, $l$.
Man setze $au \equiv v (\mathrm{mod.}\ l)$, wobei $a$ so gewählt werden kann, dass der Koeffizient von $\nu\nu$
in $v$ gleich 1 ist und für den Zahlenwert $v < \tfrac{1}{2} l$ gilt. Dann ist
\[
tv \equiv atu \equiv akl \equiv 0 (\mathrm{mod.}\ l);
\]
\noindent also setzt man
\[
\frac{tv}{l} = t',
\]
\noindent daraus folgt
\[
\varphi t' = \frac{kk\,\varphi v}{ll}
\]
\noindent und für den Zahlenwert gilt $t' < \tfrac{1}{2} t$.
\begingroup\small
\[
\begin{array}{llll}
t=1,& \varphi t = 1 & & \\
t'=v,& \varphi t' = A', & l' = \frac{A'}{(\mathrm{Div.}\,Ft')^2}, &
v' \equiv a'\,\frac{Ft'}{\mathrm{Div.}} (\mathrm{mod.}\ l'), \\
t''=\frac{t'v'}{l'}, & \varphi t'' = A'', & l' = \frac{A''}{(\mathrm{Div.}\,Ft'')^2}, & \\
t'''=\frac{t''v''}{l''}, & \varphi t''' = A''', & l''' = \frac{A'''}{(\mathrm{Div.}\,Ft''')^2}, &
\end{array}
\]
\endgroup
usw.

\Needspace{11\baselineskip}\begin{center}[2.]\end{center}
\[
\begin{gathered}
a^3+nb^3+nnc^3-3nabc = D, \\
n = \nu^3, \\
a+b\nu+c\nu\nu = \Delta,
\end{gathered}
\]
\[
\begin{array}{lll}
bb-ac = A, & C-\nu\nu A = (a-\nu b)\Delta, & CC-nAB = aD, \\
ncc-ab = B, & \nu A - B = (b-\nu c)\Delta, & nAA-BC = bD, \\
aa-nbc = C, & \nu B - C = (\nu\nu c - a)\Delta, & BB-AC = cD.
\end{array}
\]

